\section{What MuJoCo Executes}
\label{sec:mujoco-concepts}

MuJoCo is a physics engine for articulated systems, contacts, and rigid-body
dynamics. A typical program has four steps:

\begin{center}
\begin{tikzpicture}[
  node distance=8mm and 8mm,
  box/.style={draw=lirmmblue,rounded corners=1mm,thick,align=center,
              minimum height=10mm,minimum width=29mm,fill=softblue},
  arr/.style={-{Latex[length=2.2mm]},thick,lirmmblue}
]
  \node[box] (xml) {MJCF/XML\\scene};
  \node[box,right=of xml] (model) {Compile to\\\filepath{MjModel}};
  \node[box,right=of model] (data) {Create\\\filepath{MjData}};
  \node[box,right=of data] (step) {Set control and\\call \filepath{mj_step}};
  \draw[arr] (xml) -- (model);
  \draw[arr] (model) -- (data);
  \draw[arr] (data) -- (step);
\end{tikzpicture}
\end{center}

The distinction between model and data is fundamental:

\begin{description}
  \item[\filepath{mujoco.MjModel}:] the compiled, mostly fixed description of
  bodies, joints, masses, geometry, cameras, actuators, and limits.
  \item[\filepath{mujoco.MjData}:] the changing state of one simulation,
  including joint position, velocity, control input, contact information, body
  pose, and simulation time.
  \item[\filepath{mujoco.mj_step}:] advances dynamics by one model time step.
  \item[\filepath{mujoco.mj_forward}:] recomputes derived positions and
  velocities after state was changed directly, without advancing time.
\end{description}

An XML edit requires the model to be reloaded. In contrast, moving a free object
at reset normally modifies \filepath{MjData} and then calls
\filepath{mj_forward}.

\section{Run the Minimal Example}

The repository contains a small example independent of the microwave task:

\begin{itemize}
  \item \filepath{dp_mujoco/examples/minimal_scene.xml};
  \item \filepath{dp_mujoco/examples/mujoco_basics.py}.
\end{itemize}

It contains a floor, light, one actuated hinge, a free cube, a site, and a
camera. Run it without graphics first:

\begin{lstlisting}[language=bash]
python dp_mujoco/examples/mujoco_basics.py
\end{lstlisting}

The output reports the simulation time step and the dimensions \filepath{nq},
\filepath{nv}, and \filepath{nu}, then prints the changing joint position and
control target. Open the same example interactively with:

\begin{lstlisting}[language=bash]
python dp_mujoco/examples/mujoco_basics.py --viewer --duration 10
\end{lstlisting}

The blue arm follows a sinusoidal target, the green cube falls under gravity,
and the red site follows the end of the arm. This one scene demonstrates
actuation, free-body dynamics, contact, and a kinematic reference frame.

\section{Read a Small MJCF Scene}

MJCF describes a hierarchy. Child bodies move with their parent, while joints
define the motion allowed between them. The important elements are:

\begin{table}[H]
  \centering
  \begin{tabularx}{\textwidth}{@{}p{0.20\textwidth}X@{}}
    \toprule
    \textbf{Element} & \textbf{Purpose} \\
    \midrule
    \filepath{body} & A rigid frame that can contain joints, geometry, sites, and child bodies. \\
    \filepath{joint} & One degree of freedom such as a hinge or slide. \\
    \filepath{freejoint} & Six free degrees of freedom, represented by 3D position and a quaternion. \\
    \filepath{geom} & Visual and/or collision geometry with size, mass, friction, and material. \\
    \filepath{site} & A massless reference frame used for a TCP, sensor location, or target. \\
    \filepath{camera} & A fixed or body-mounted viewpoint used by the viewer or renderer. \\
    \filepath{actuator} & Converts a value in \filepath{data.ctrl} into force, velocity, or position control. \\
    \bottomrule
  \end{tabularx}
\end{table}

The core of the tutorial scene is deliberately small:

\begin{lstlisting}[language=XML]
<body name="arm" pos="0 0 0.08">
  <joint name="arm_hinge" type="hinge" axis="0 1 0"
         range="-1.57 1.57"/>
  <geom type="capsule" fromto="0 0 0 0 0 0.35"
        size="0.025" mass="0.5"/>
  <site name="tool_site" pos="0 0 0.35" size="0.015"/>
</body>

<actuator>
  <position name="arm_position" joint="arm_hinge" kp="20"/>
</actuator>
\end{lstlisting}

The actuator type defines the meaning of \filepath{data.ctrl}. Here the value
is a desired hinge angle. A motor actuator would interpret it differently, so
never assume that \filepath{ctrl} always contains joint positions.

MuJoCo uses half-extents for a box \filepath{geom}: a box with
\filepath{size="0.04 0.04 0.04"} measures $8\times8\times8$~cm. Incorrectly
treating these values as full dimensions is a common source of collision errors.

\section{Understand the Python Simulation Loop}

The essential Python loop is:

\begin{lstlisting}[language=Python]
model = mujoco.MjModel.from_xml_path("minimal_scene.xml")
data = mujoco.MjData(model)

actuator_id = mujoco.mj_name2id(
    model, mujoco.mjtObj.mjOBJ_ACTUATOR, "arm_position"
)

while data.time < 4.0:
    data.ctrl[actuator_id] = target_angle
    mujoco.mj_step(model, data)
\end{lstlisting}

The most frequently inspected arrays are:

\begin{description}
  \item[\filepath{data.qpos}:] generalized positions. A hinge uses one value;
  a free joint uses seven values: three for position and four for orientation.
  \item[\filepath{data.qvel}:] generalized velocities. A free joint uses six
  values, so \filepath{nq} and \filepath{nv} need not be equal.
  \item[\filepath{data.ctrl}:] one input per actuator; its length is
  \filepath{model.nu}.
  \item[\filepath{data.site_xpos/site_xmat}:] world position and orientation of
  sites after forward kinematics or a simulation step.
\end{description}

Use names instead of hard-coded indices when possible. The project obtains the
TCP site in the same way:

\begin{lstlisting}[language=Python]
site_id = mujoco.mj_name2id(
    model, mujoco.mjtObj.mjOBJ_SITE, "grasp_site"
)
tcp_position = data.site_xpos[site_id].copy()
\end{lstlisting}

\section{Open the Project Scene}
\label{sec:mujoco-scene}

The active task scene is:

\begin{center}
\filepath{dp_mujoco/models/universal_robots_ur10e/scene_microwave_camera.xml}
\end{center}

The scene includes \filepath{ur10e_with_gripper.xml}; this keeps robot links,
joints, gripper, actuators, wrist camera, and \filepath{grasp_site} separate
from the microwave, table, objects, target boxes, lighting, and global camera.
The Python files around the XML are:

\begin{table}[H]
  \centering
  \begin{tabularx}{\textwidth}{@{}p{0.42\textwidth}X@{}}
    \toprule
    \textbf{Path} & \textbf{Responsibility} \\
    \midrule
    \filepath{dp_mujoco/env/mujoco_env.py} & Loads the scene and exposes state, rendering, reset, and commands. \\
    \filepath{dp_mujoco/env/scene_utils.py} & Changes object pose during teleoperation resets. \\
    \filepath{dp_mujoco/utils/scene_utils.py} & Applies the matching reset distribution during policy execution. \\
    \filepath{dp_mujoco/teleop/test_UR10e_touch.py} & Runs simulated teleoperation and recording. \\
    \filepath{dp_mujoco/policy_exec/run_diffusion_mujoco.py} & Executes a trained checkpoint in MuJoCo. \\
    \bottomrule
  \end{tabularx}
\end{table}

Open the XML without ROS or the Touch:

\begin{lstlisting}[language=bash]
python -m mujoco.viewer \
  --mjcf=dp_mujoco/models/universal_robots_ur10e/scene_microwave_camera.xml
\end{lstlisting}

Inspect the body hierarchy, joint limits, contact geometry, and initial object
poses. The rendered mesh may look correct even when its collision geometry is
wrong, so enable contact or geometry visualization when debugging.

\section{Separate Visual and Collision Geometry}

A detailed mesh is useful for appearance but often unnecessarily difficult for
contact. In the project scene, the microwave OBJ mesh has collision disabled and
simple transparent boxes represent its floor, walls, and back. This makes
contact faster and easier to inspect.

When adding an object, begin with a primitive:

\begin{lstlisting}[language=XML]
<body name="practice_cube" pos="0.60 0.00 0.48">
  <freejoint/>
  <geom type="box" size="0.025 0.025 0.025"
        mass="0.10" rgba="0.15 0.55 0.90 1"/>
</body>
\end{lstlisting}

Reload the XML and let the object settle. If it falls through, jumps, or becomes
stuck, check units, initial overlap, collision dimensions, mass, friction, and
contact settings before modifying the robot controller.

\section{Render the Policy Cameras}

The project uses \filepath{top_table} as the global camera and
\filepath{wrist_cam} as the eye-in-hand camera. A renderer must be updated from
the latest \filepath{MjData} before an image is read:

\begin{lstlisting}[language=Python]
renderer = mujoco.Renderer(model, height=480, width=640)
renderer.update_scene(data, camera="top_table")
rgb_image = renderer.render()
\end{lstlisting}

\begin{figure}[H]
  \centering
  \includegraphics[width=0.86\textwidth]{camera_view.png}
  \caption{Project camera observations: wrist view on the left and global view on the right.}
  \label{fig:sim-camera-example}
\end{figure}

The teleoperation display uses $640\times480$ images, while its recorder resizes
them to $84\times84$ for the simulation dataset. Camera name, ordering,
orientation, and final image size must remain identical during policy execution.

\section{Reset and Randomize an Object}

A free joint stores object position followed by a quaternion in
\filepath{data.qpos}. The project reset utility finds the free-joint address,
writes the new pose, clears its six velocity values, and calls
\filepath{mj_forward}:

\begin{lstlisting}[language=Python]
qpos_adr = model.jnt_qposadr[joint_id]
qvel_adr = model.jnt_dofadr[joint_id]

data.qpos[qpos_adr:qpos_adr + 3] = position
data.qpos[qpos_adr + 3:qpos_adr + 7] = quat_wxyz
data.qvel[qvel_adr:qvel_adr + 6] = 0.0
mujoco.mj_forward(model, data)
\end{lstlisting}

The reset distribution used for collection and policy execution must match.
The current repository has one randomization module for each path; update and
test both when the pose range changes. Run many resets before recording and
confirm that every object pose is collision-free, visible, reachable, and part
of the intended task.

\begin{checkpointbox}[MuJoCo environment ready]
The minimal example runs, the project XML opens, the robot and free objects
behave correctly under gravity and contact, both cameras render useful views,
and repeated reset samples remain valid. Only then connect the Touch and test
the complete teleoperation loop.
\end{checkpointbox}
